\documentclass[aps,prd,twocolumn,superscriptaddress,nofootinbib]{revtex4-2}

\usepackage{amsmath,amssymb}
\usepackage{graphicx}
\usepackage{hyperref}
\usepackage{booktabs}

\begin{document}

\title{Probing geometry in causal set quantum gravity:\\spectral dimension fails, vacuum entanglement succeeds}

\author{Matt Loftus}
\email{matthew.a.loftus@gmail.com}
\affiliation{Independent researcher}

\date{\today}

\begin{abstract}
We study how to detect emergent geometry on causal sets sampled from the Benincasa-Dowker partition function. We prove that the $d_s = 2$ crossing of the spectral dimension is generic to all finite connected graphs (by the intermediate value theorem), show that the peak $d_s$ is determined by link density (random graphs match causal sets), and demonstrate that the CDT plateau grows with system size while the causet plateau shrinks. We then show that Sorkin-Johnston vacuum entanglement entropy provides a genuinely geometric probe: it scales as $S \sim \ln N$ consistent with 1+1D conformal field theory, and drops by a factor of $3.4$ across the Benincasa-Dowker phase transition. For spatially motivated partitions, the continuum phase satisfies the holographic monogamy inequality $I_3 \leq 0$ significantly above chance ($p < 10^{-7}$), though this result is partition-dependent. The SJ entanglement entropy distinguishes the phases of causal set quantum gravity in cases where connectivity-based observables cannot.
\end{abstract}

\maketitle

\section{Introduction}

A fundamental challenge in discrete approaches to quantum gravity is identifying whether a given discrete structure approximates a smooth spacetime. The spectral dimension $d_s$---defined via the return probability of a diffusion process---has been the primary tool, with dimensional reduction from $d_s = 4$ to $d_s \approx 2$ at short scales reported in at least eight independent approaches~\cite{carlip2017}. However, the reliability of this observable on causal sets~\cite{bombelli1987,surya2019} has been questioned~\cite{eichhorn2014,eichhorn2019}.

In this paper, we present evidence that link-graph spectral dimension on causal sets is determined by connectivity, not geometry (Sec.~II--III). We then introduce a replacement observable---the Sorkin-Johnston (SJ) vacuum entanglement entropy~\cite{johnston2009,sorkin2012}---that genuinely probes causal geometry (Sec.~IV--V). The combined narrative moves from a negative result (what doesn't work) to a positive one (what does, and why it's holographic).

\section{Spectral dimension: the crossing theorem}

The spectral dimension is defined via the normalized graph Laplacian $L = I - D^{-1/2}AD^{-1/2}$ with eigenvalues $\{\lambda_k\}$. The return probability $P(\sigma) = n^{-1}\sum_k e^{-\lambda_k\sigma}$ gives $d_s(\sigma) = -2\,d\ln P/d\ln\sigma$.

\textbf{Proposition.} \emph{For any finite connected graph, $d_s(\sigma)$ is smooth on $(0,\infty)$ with $d_s \to 0$ as $\sigma \to 0^+$ and $\sigma \to \infty$. If $d_s^{\max} > c$, the IVT gives $\sigma_1 < \sigma_2$ with $d_s(\sigma_1) = d_s(\sigma_2) = c$.}

The proof follows from $P > 0$ (by connectedness), the vanishing of $\sigma P'$ at both limits, and continuity. The value $c = 2$ plays no special role: the crossing is topological, not physical.

\section{Spectral dimension fails on causal sets}

We compare sprinkled 2D causal sets, 2D CDT, random DAGs, and Erd\H{o}s-R\'enyi graphs at $N = 100$--$800$.

\textbf{Result 1:} All graph types cross $d_s = 2$ (Table~\ref{tab:crossing}, Fig.~\ref{fig:combined}a). Random graphs produce peak $d_s$ values indistinguishable from those of causal sets.

\begin{figure*}[t]
\includegraphics[width=0.48\textwidth]{fig1_combined.pdf}
\hfill
\includegraphics[width=0.48\textwidth]{fig2_entanglement.pdf}
\caption{\textbf{Left:} (a) Peak spectral dimension at $N = 500$. Random graphs (red) match causal sets (blue); only CDT and lattices recover the correct dimension. (b) Holographic monogamy test: the continuum phase satisfies $I_3 \leq 0$ in 97\% of tripartitions ($p = 8.7 \times 10^{-8}$ above the 52\% random baseline), while the crystalline phase satisfies it in only 3\%. \textbf{Right:} (a) SJ entanglement entropy $S(N/2)/\ln N$ converges to $\approx 1.0$ (CFT-like scaling). (b) Entanglement profiles in the continuum (green) and crystalline (red) phases, showing a $3.4\times$ suppression.}
\label{fig:combined}
\end{figure*}

\begin{table}[h]
\centering
\begin{tabular}{lc}
\toprule
Graph type & Peak $d_s$ \\
\midrule
2D CDT & 2.54 \\
Sprinkled 2D causet & 5.21 \\
Random DAG & 4.93 \\
Erd\H{o}s-R\'enyi & 4.95 \\
2D lattice & 2.29 \\
\bottomrule
\end{tabular}
\caption{Peak spectral dimension at $N = 500$. Random graphs produce values comparable to causal sets.}
\label{tab:crossing}
\end{table}

\textbf{Result 2:} CDT has a growing plateau while the causet plateau shrinks (Table~\ref{tab:plateau}). CDT's peak $d_s$ remains stable near 2.5 (correctly reflecting 2D); the causet peak grows to 5.66. The plateau width is defined as the fraction of $\sigma$ values with $|d_s - 2| < 0.3$ (results are qualitatively unchanged for thresholds 0.2 and 0.5).

\begin{table}[h]
\centering
\begin{tabular}{lcccc}
\toprule
& \multicolumn{2}{c}{CDT} & \multicolumn{2}{c}{Causet} \\
\cmidrule(lr){2-3} \cmidrule(lr){4-5}
$N$ & Peak & Plateau & Peak & Plateau \\
\midrule
100 & 2.49 & 17\% & 3.64 & 8\% \\
400 & 2.50 & 36\% & 4.98 & 7\% \\
800 & 2.59 & 47\% & 5.66 & 6\% \\
\bottomrule
\end{tabular}
\caption{CDT plateau grows with $N$; causet plateau shrinks. Both in 2D.}
\label{tab:plateau}
\end{table}

\textbf{Result 3:} Random graphs with matched link density produce the same peak $d_s$ as causal sets ($3.0$ at $L/N = 2.6$; $4.0$ at $L/N = 10.6$). This holds for both the link-graph Laplacian and the Benincasa-Dowker d'Alembertian~\cite{benincasa2010}. The problem is structural, not operational.

\textbf{Reinterpretation:} Eichhorn and Mizera~\cite{eichhorn2014} found that $d_s$ \emph{increases} at small scales on causal sets---opposite to CDT. Our finding explains this: the walker probes the high-connectivity local neighborhood, and $d_s$ tracks connectivity. Subsequent fixes (spatial hypersurfaces~\cite{eichhorn2019}, continuum d'Alembertian~\cite{belenchia2016}) all work by taming the scale-dependent connectivity.

\section{SJ vacuum entanglement: construction}

The Sorkin-Johnston vacuum for a free massless scalar on a causal set $\mathcal{C}$ with causal matrix $C$ is constructed from the Pauli-Jordan function $\Delta_{ij} = (2/N)(C_{ji} - C_{ij})$, a real antisymmetric matrix. The Hermitian matrix $i\Delta$ has real eigenvalues $\{\lambda_k\}$ in $\pm$ pairs, and the Wightman function is the positive-frequency part:
\begin{equation}
W = \sum_{\lambda_k > 0} \lambda_k \, |\mathbf{v}_k\rangle\langle\mathbf{v}_k|
\end{equation}
\cite{johnston2009}. The $2/N$ normalization ensures $W$ eigenvalues lie in $[0,1]$. The entanglement entropy of sub-region $A$ is $S(A) = -\mathrm{Tr}[W_A\ln W_A + (I-W_A)\ln(I-W_A)]$~\cite{sorkin2012}.

The failure of spectral dimension raises a concrete question: what observable \emph{does} probe geometry on causal sets? We propose the Sorkin-Johnston vacuum entanglement entropy and test it across the same BD transition where spectral dimension fails.

For the BD transition measurements, we sample 2-orders from the partition function $Z = \sum_\mathcal{C} e^{-\beta S[\mathcal{C}]/\hbar}$ using Metropolis coordinate-swap moves on the space of $N$-element 2-orders~\cite{surya2012}, with the non-local action at $\epsilon = 0.12$~\cite{glaser2018}. We use $5 \times 10^4$ MCMC steps with $2.5 \times 10^4$ thermalization, and parallel tempering with 8 replicas for large $\beta$.

For the monogamy test, we partition elements into thirds using the $V$-coordinate rank of the generating 2-order permutation. For BD-sampled causets (which lose the permutation after MCMC), we use an index-based tripartition (first, middle, and last thirds by label).

\section{SJ entanglement probes geometry}

\subsection{Logarithmic scaling}

For random 2-orders ($\beta = 0$), $S(N/2)/\ln N$ converges to $\approx 1.0$ at small $N$ (Table~\ref{tab:scaling}), consistent with 1+1D CFT scaling $S = (c/3)\ln N$~\cite{calabrese2004} with $c \approx 3$. The $c > 1$ value reflects contributions from untruncated near-zero modes of $i\Delta$~\cite{sorkin_yazdi2018}; the relative comparison between phases is robust to normalization.

\begin{table}[h]
\centering
\begin{tabular}{cccc}
\toprule
$N$ & $S(N/2)$ & $S/\ln N$ \\
\midrule
20 & $2.56 \pm 0.15$ & $0.85 \pm 0.05$ \\
30 & $3.15 \pm 0.12$ & $0.93 \pm 0.04$ \\
40 & $3.58 \pm 0.10$ & $0.97 \pm 0.03$ \\
50 & $3.92 \pm 0.09$ & $1.00 \pm 0.02$ \\
\bottomrule
\end{tabular}
\caption{Entanglement entropy scaling in the continuum phase at small $N$. $S/\ln N \to 1.0$ at $N = 50$.}
\label{tab:scaling}
\end{table}

\textbf{Large-$N$ caveat.} At larger system sizes ($N = 100$--$500$), the effective central charge $c_{\mathrm{eff}} = 3S(N/2)/\ln N$ continues to grow: $c_{\mathrm{eff}} = 3.0$ ($N = 50$), $3.3$ ($N = 100$), $3.7$ ($N = 200$), $4.1$ ($N = 500$). This growth persists in MCMC-sampled continuum-phase 2-orders (at $\beta = 0.7\beta_c$), indicating that it is intrinsic to the SJ vacuum on 2-orders with the $2/N$ normalization, not an artifact of non-manifold-like configurations. The entropy scaling is thus super-logarithmic at large $N$, not converging to $c = 1$ as expected for a free scalar in continuum 2D CFT. The phase comparison ($3.4\times$ drop across the BD transition) remains valid as a relative diagnostic, but the absolute value of $c$ on 2-orders should not be compared directly with continuum CFT predictions.

\subsection{Phase transition}

Across the BD transition ($N = 40$, $\epsilon = 0.12$), $S(N/2)$ drops from $3.71 \pm 0.10$ (continuum) to $1.10$ (crystalline)---a factor of 3.4.

\subsection{Holographic monogamy}

The tripartite information $I_3(A{:}B{:}C) = I(A{:}B) + I(A{:}C) - I(A{:}BC)$ satisfies $I_3 \leq 0$ (monogamy) in holographic theories~\cite{hayden2013}. We find (Table~\ref{tab:monogamy}):

\begin{table}[h]
\centering
\begin{tabular}{lcc}
\toprule
Phase & $\langle I_3 \rangle$ & $P(I_3 \leq 0)$ \\
\midrule
Continuum & $-0.10 \pm 0.04$ & 97\% \\
Crystalline & $-0.00 \pm 0.00$ & 3\% \\
Random Gaussian & $-0.00 \pm 0.02$ & 52\% \\
\bottomrule
\end{tabular}
\caption{Monogamy test. The continuum phase is significantly more monogamous than random (97\% vs 52\%); the crystalline phase is significantly less (3\%).}
\label{tab:monogamy}
\end{table}

The continuum phase satisfies monogamy far above chance (97\% vs 52\% for random Gaussian states; $p = 8.7 \times 10^{-8}$ by a one-sided binomial test against the 52\% null rate, $n = 30$), while the crystalline phase violates it ($p = 9.2 \times 10^{-9}$). The Hayden-Headrick-Maloney inequality~\cite{hayden2013} was proven for large-$c$ holographic CFTs; our $c \approx 3$ is outside this regime, so the comparison is suggestive rather than rigorous. Nevertheless, the sharp contrast between phases demonstrates that monogamy is a property of causal structure, not the Gaussian formalism.

\textbf{Partition dependence.} An important caveat: the monogamy result depends on the partition scheme. The 97\% rate is obtained using the $V$-coordinate of the generating 2-order, which provides a geometrically natural spatial partition. When we instead partition by element index (which is available for arbitrary causal sets without a 2-order representation), the monogamy rate in the continuum phase drops to $\sim$50\%---indistinguishable from chance. Furthermore, random DAGs with index-based partitions satisfy monogamy in $\sim$100\% of cases, producing false positives. This means $I_3$ is a reliable diagnostic only when a geometrically meaningful partition is available. The phase comparison in Table~\ref{tab:monogamy} remains valid because both phases use the same partition scheme, but $I_3$ should not be interpreted as an intrinsic property of the causal set without specifying the partition.

\section{Discussion}

The combined picture:
\begin{enumerate}
\item \textbf{Spectral dimension fails} because it tracks connectivity (link density), not geometry. This is proven by the crossing theorem (any graph crosses $d_s = 2$), the scaling data (CDT plateau grows, causet shrinks), the link-density control (random graphs match), and the d'Alembertian test (same failure with the BD operator).
\item \textbf{SJ entanglement succeeds} because it captures quantum correlations induced by causal structure. It produces CFT-like $\ln N$ scaling, distinguishes the BD phases by a factor of 3.4, and satisfies holographic monogamy in the continuum phase.
\item \textbf{Monogamy as supporting evidence.} For spatially motivated partitions, the continuum phase satisfies $I_3 \leq 0$ in 97\% of cases (vs.\ 52\% for random Gaussian states). However, this result is partition-dependent and should be interpreted as supporting evidence for geometric entanglement structure, not as an intrinsic property of the causal set.
\end{enumerate}

We note that the SJ entanglement entropy correlates strongly with the BD action across the transition ($r > 0.98$), but the proportionality constant depends on $N$ and $\epsilon$ and is not universal.

The practical recommendation is that causal set studies seeking to detect emergent geometry should use the SJ vacuum entanglement entropy rather than link-graph spectral dimension.

\section{Conclusion}

We have shown that link-graph spectral dimension cannot probe geometry on causal sets (it is determined by connectivity), while the Sorkin-Johnston vacuum entanglement entropy can (it exhibits CFT-like logarithmic scaling and sharply distinguishes the BD phases). For geometrically motivated partitions, the continuum phase additionally satisfies holographic monogamy, though this result is partition-dependent. The SJ entanglement entropy is a more reliable probe of emergent geometry than connectivity-based observables in discrete quantum gravity.

\begin{acknowledgments}
This work was conducted with assistance from Claude (Anthropic).
\end{acknowledgments}

\begin{thebibliography}{99}

\bibitem{carlip2017}
S.~Carlip,
``Dimension and dimensional reduction in quantum gravity,''
Class.\ Quant.\ Grav.\ \textbf{34}, 193001 (2017),
\href{https://arxiv.org/abs/1705.05417}{arXiv:1705.05417}.

\bibitem{eichhorn2014}
A.~Eichhorn and S.~Mizera,
``Spectral dimension in causal set quantum gravity,''
Class.\ Quant.\ Grav.\ \textbf{31}, 125007 (2014),
\href{https://arxiv.org/abs/1311.2530}{arXiv:1311.2530}.

\bibitem{eichhorn2019}
A.~Eichhorn, S.~Surya, and F.~Versteegen,
``Spectral dimension on spatial hypersurfaces in causal set quantum gravity,''
Class.\ Quant.\ Grav.\ \textbf{36}, 205013 (2019),
\href{https://arxiv.org/abs/1905.13498}{arXiv:1905.13498}.

\bibitem{johnston2009}
S.~Johnston,
``Feynman propagator for a free scalar field on a causal set,''
Phys.\ Rev.\ Lett.\ \textbf{103}, 180401 (2009),
\href{https://arxiv.org/abs/0909.0944}{arXiv:0909.0944}.

\bibitem{sorkin2012}
R.~D.~Sorkin,
``Expressing entropy globally in terms of (4D) field-correlations,''
J.\ Phys.\ Conf.\ Ser.\ \textbf{484}, 012004 (2014),
\href{https://arxiv.org/abs/1205.2953}{arXiv:1205.2953}.

\bibitem{benincasa2010}
D.~M.~T.~Benincasa and F.~Dowker,
``The scalar curvature of a causal set,''
Phys.\ Rev.\ Lett.\ \textbf{104}, 181301 (2010),
\href{https://arxiv.org/abs/1001.2725}{arXiv:1001.2725}.

\bibitem{belenchia2016}
A.~Belenchia, D.~M.~T.~Benincasa, A.~Marcian\`o, and L.~Modesto,
``Spectral dimension from nonlocal dynamics on causal sets,''
Phys.\ Rev.\ D \textbf{93}, 044017 (2016),
\href{https://arxiv.org/abs/1507.00330}{arXiv:1507.00330}.

\bibitem{hayden2013}
P.~Hayden, M.~Headrick, and A.~Maloney,
``Holographic mutual information is monogamous,''
Phys.\ Rev.\ D \textbf{87}, 046003 (2013),
\href{https://arxiv.org/abs/1107.2940}{arXiv:1107.2940}.

\bibitem{sorkin_yazdi2018}
R.~D.~Sorkin and Y.~K.~Yazdi,
``Entanglement entropy in causal set theory,''
Class.\ Quant.\ Grav.\ \textbf{35}, 074004 (2018),
\href{https://arxiv.org/abs/1611.10281}{arXiv:1611.10281}.

\bibitem{glaser2018}
L.~Glaser, D.~O'Connor, and S.~Surya,
``Finite size scaling in 2d causal set quantum gravity,''
Class.\ Quant.\ Grav.\ \textbf{35}, 045006 (2018),
\href{https://arxiv.org/abs/1706.06432}{arXiv:1706.06432}.

\bibitem{bombelli1987}
L.~Bombelli, J.~Lee, D.~Meyer, and R.~D.~Sorkin,
``Space-time as a causal set,''
Phys.\ Rev.\ Lett.\ \textbf{59}, 521 (1987).

\bibitem{surya2019}
S.~Surya,
``The causal set approach to quantum gravity,''
Living Rev.\ Rel.\ \textbf{22}, 5 (2019),
\href{https://arxiv.org/abs/1903.11544}{arXiv:1903.11544}.

\bibitem{surya2012}
S.~Surya,
``Evidence for the continuum in 2D causal set quantum gravity,''
Class.\ Quant.\ Grav.\ \textbf{29}, 132001 (2012),
\href{https://arxiv.org/abs/1110.6244}{arXiv:1110.6244}.

\bibitem{calabrese2004}
P.~Calabrese and J.~Cardy,
``Entanglement entropy and quantum field theory,''
J.\ Stat.\ Mech.\ \textbf{0406}, P06002 (2004),
\href{https://arxiv.org/abs/hep-th/0405152}{arXiv:hep-th/0405152}.

\end{thebibliography}

\end{document}
